\documentclass[10pt,a4paper]{amsart}
\usepackage[margin=19mm]{geometry}
\usepackage{amsmath,amssymb,mathtools}
\usepackage[scr=rsfs,cal=euler]{mathalfa}
\usepackage{xfrac}
\usepackage[unicode,hidelinks,bookmarks=false]{hyperref}
\newcommand{\ie}{i.e.\ }
\newcommand{\cf}{cf.\ }
\newcommand{\resp}{respectively\ }
\newcommand{\Subsection}{\subsection}
\providecommand{\nobreakdash}{\nobreak\hbox{-}}
\setlength{\emergencystretch}{2em}
\multlinegap0pt
\usepackage{bm}
\usepackage{bbm}
\usepackage{dsfont}
\usepackage{xspace}
\usepackage{cancel}
\usepackage{mathtools}
\usepackage{amsthm}
\makeatletter
\@ifundefined{lem}{\newtheorem{lem}{Lemma}}{}
\@ifundefined{prop}{\newtheorem{prop}{Proposition}}{}
\makeatother
\newcommand{\bx}{\mathbf{x}}
\newcommand{\by}{\mathbf{y}}
\newcommand{\xconst}{X_{s}}
\title[Cellular automata can really solve the parity problem]{Cellular automata can really solve the parity problem}
\author{Barbara Wolnik, Anna Nenca, Pedro Paulo Balbi, Bernard De Baets}
\date{Source: arXiv:2501.08684v2 (CC BY 4.0)}
\begin{document}
\maketitle

\noindent\emph{Source and licence.} Cellular automata can really solve the parity problem. Version arXiv:2501.08684v2 (CC BY 4.0), distributed under the Creative Commons Attribution 4.0 International licence, as recorded in the arXiv licence metadata for this exact version and shown on the abstract page https://arxiv.org/abs/2501.08684v2. Full rights notice: licenses/arxiv-2501.08684-CC-BY-4.0.txt.\par\smallskip

\noindent\textit{Editorial scope.} Counted unit: the Prop whose source label is \texttt{lem:snakes} in the article (source0.tex lines 960--963, source label \texttt{lem:snakes}), delivered together with the article's own complete original proof of it (source0.tex lines 964--999, one \texttt{proof} environment of 36 lines). No mathematics is written, repaired or re-derived: statement and proof are the article's own text line for line. Required context retained uncounted: thm:5-12 (lines 879--881); lem:alter (lines 936--938); lem:D56 (lines 944--946); lem:longer (lines 951--953). Every definition, display and label of those lines is retained unchanged. Omitted from this package and not redistributed: 1--878, 882--935, 939--943, 947--950, 954--959, 1000--1048. Nothing inside a retained excerpt is omitted or altered: every retained body is delivered exactly as the article's lines are written, so no omission receipt of any kind accompanies this package. Citations: the retained excerpts carry no citation command at all, so no bibliography entry of the article is transcribed and no reference is redistributed. Reference adaptation: none was needed and none was made; every reference inside the delivered unit resolves inside it, so no unresolved reference or citation is produced. Printed slips kept literal: no printed word, symbol, hypothesis or number of the counted statement or of its proof was altered. Layout and class: the article's own class is \texttt{sn-jnl} with packages including amsthm, graphicx, multirow, amsmath, amssymb, amsfonts, mathrsfs, appendix, xcolor, textcomp, manyfoot, booktabs, algorithm, algorithmicx; the wrapper is the lane's pinned \texttt{amsart} editorial wrapper at 10pt on A4, which restates the theorem-like environments the retained text needs and adds the article's own macro definitions that the retained text needs (\texttt{\textbackslash bx}, \texttt{\textbackslash by}, \texttt{\textbackslash xconst}), with extra wrapper packages (bm, bbm, dsfont, xspace, cancel, mathtools). The delivered numbering is the unit's own. Assets: no retained span contains a figure environment, a graphics-inclusion command, a PDF-page inclusion, a table or a TikZ picture; no image file is copied into this package, the package declares no asset dependency and no image file is redistributed with it. Licence: version arXiv:2501.08684v2 of this article is distributed under the Creative Commons Attribution 4.0 International licence, as recorded in the arXiv licence metadata for this exact version and shown on the abstract page https://arxiv.org/abs/2501.08684v2. A full rights notice is delivered at licenses/arxiv-2501.08684-CC-BY-4.0.txt. Characters: every retained excerpt is pure ASCII, so the delivered engine, XeLaTeX with its default Latin Modern text font, reports no missing character.
\begin{prop}\label{thm:5-12}
    Let $\bx\in \xconst$. Then $\bx$ does not contain $D_{5,6}^r$, $D_{7,8}$, $D_{9,10}^r$ or $D_{9,12}^r$ and the action of $T_{1,2}$ and $T_{3,4}$ cannot lead to the merging of two blocks of 1s. Moreover, if $b=x_{i}x_{i+1}\ldots x_{i+2k+1}$ is an ordered block in $\bx$ and $\by=F(\bx)$, then there is an ordered block in $\by$ containing $y_{i}y_{i+1}\ldots y_{i+2k+1}$.
\end{prop}

\begin{lem}\label{lem:alter}
If $\bx\in \xconst$, then $\bx$ does not contain $010101$.
\end{lem}

\begin{lem}\label{lem:D56}
Let $\bx\in \xconst$ and let $b$ be an ordered block in $\bx$. Then $b$ does not contain $001100$. Also, if $b$ is maximal, then it does not end with $0011$.
\end{lem}

\begin{lem}\label{lem:longer}
Let $\bx\in \xconst$ and let $b$ be any maximal ordered block in $\bx$. If $b$ ends with $11$, then $F(\bx)$ contains an ordered block longer than $b$.
\end{lem}

\begin{prop}
\label{lem:snakes}
If $\bx\in \xconst$, then $\bx$ does not contain any ordered block.
\end{prop}

\begin{proof}
First of all, let us recall that if $\bx\in\xconst$, then, according to Proposition~\ref{thm:5-12}, the ordered blocks in $\bx$ do not get shortened as time evolves.

Let us assume that some configurations from $\xconst$ contain an ordered block, from which we choose those that have the shortest length, and then, from them, we choose  those that contain the longest ordered block (whose length we denote by $l_0$). Let us denote the set of all such configurations by $B$. 
Let $\bx\in B$ and, without loss of generality, assume that $b=x_{0}x_1\ldots x_{2k+1}$ is the longest ordered block in $\bx$ (so, it is maximal). 

First we show that $b$ cannot end with $11$. 
Indeed, according to Lemma~\ref{lem:longer}, if $b$ ends with $11$, there should be a longer ordered block than $b$ in $F(\bx)$, contradicting our choice of $\bx$ and $b$.

Next, we analyse whether $b$ might end with $00$. Thus, suppose that $x_{2k}x_{2k+1}=00$ and let $x_{2m}x_{2m+1}$ be the last two-symbol block in $b$ that is not equal to $00$. If $x_{2m}x_{2m+1}=11$, then also  $x_{2m-1}=1$ (according to Lemma~\ref{lem:D56}), which means that the end of $b$ would look as follows: 
\[
b = \ ?\ldots?111\underbrace{00\ldots 00}_{2(k-m)}\ ;
\]
however, due to $T_{1,2}$, after $k-m$ time steps the ordered block $b$ would evolve into $b'$ in $F^{k-m}(\bx)$, having at least the same length as $b$ (so, $F^{k-m}(\bx)\in B$), but ending with a block of 1s, which, as already shown, would be a contradiction. Therefore, $b$ also cannot end with a block of at least two 0s which is preceded by at least two 1s.

So far, we have shown that if $B$ is non-empty, then every ordered block of maximum length contained in $\bx\in B$ ends with a block of 0s preceded by $01$. Now let us take such $\bx\in B$ in which an ordered block of maximum length $b=x_{0}x_1\ldots x_{2k+1}$ has the longest block of 0s at the end. This means that the end of $b$ would look as follows: 
\[
b=\ ?\ldots?01\underbrace{00\ldots 00}_{2(k-m)}\ .
\]
Note that $x_{2m}x_{2m+1}\ldots x_{2k}x_{2k+1}=0100\ldots 00$ cannot be preceded by $111$, as then there would be an ordered block $b'$ in $F(\bx)$, at least as long as $b$, and ending with the longer block of 0s:
\[
\phantom{ff}\bx=\ ?\ldots?11101\underbrace{00\ldots 00}_{2(k-m)}? \ , 
\]
\[
F(\bx)=\ ?\ldots?10100\underbrace{00\ldots 00}_{2(k-m)}? \ ,
\]
contradicting our choice of $\bx$ and $b$.

All that remains is to check all other possibilities for the part of $\bx$ preceding $x_{2m}x_{2m+1}\ldots x_{2k}x_{2k+1}=0100\ldots 00$, which leads to four possibilities:
\begin{itemize}\itemindent=15pt
    \item[Case 1:] If $x_{2m-1}=0$, due to $T_{3,4}$, after $k-m$ time steps the ordered block $b$ would evolve into $b'$ ending with $11$, which, as we already know, cannot happen.
    \item[Case 2:] If $x_{2m-3}x_{2m-2}x_{2m-1}=001$, then $\bx$ contains $D_{7,8}$, which is forbidden according to Proposition~\ref{thm:5-12}.
    \item[Case 3:] If $x_{2m-3}x_{2m-2}x_{2m-1}=011$, then due to $T_{5,6}$, $b$ would evolve into $b'$ in $F(\bx)$ with at least the same length as $b$, but ending with $00b$ and we are in Case 1.
    \item[Case 4:] If $x_{2m-3}x_{2m-2}x_{2m-1}=101$, then $x_{2m-4}=1$ (due to Lemma~\ref{lem:alter}) and $x_{2m-5}=1$ (otherwise, $\bx$ would contain $D_{5,6}^{r}$). But then $F(\bx)$ would contain $00b$, which leads again to Case 1.
\end{itemize}
\end{proof}

\end{document}
